% SGA 3, Expose VI_B, Loop 1 English reconstruction.
% Scope: Proposition 6.2, Definition 6.2.1, and Proposition 6.2.2
% in full, including source editor's notes 51--54.
% Authority: Polo--Gille Exp6B-13oct24.pdf, component-local pages
% 39--40 (combined-reader pages 380--381; displayed re-edition
% headers 369--370). Theorem 6.2.3 begins on local page 40 and is
% excluded.
% OCR and Reinhold English Markdown are comparison witnesses only.
% This unit contains no diagrams.

\medskip
\phantomsection\label{prop:sga3-VIB-6.2}%
\noindent\textbf{命题 6.2.---}%
\footnote{\label{note:sga3-VIB-51}编者注: 我们按编者\hyperref[note:sga3-VIB-49]{注 49}
所述修改了陈述.}
\emph{设 \(G\) 为 \(S\)-群. 对函子 \(G\) 的子函子 \(T\),
考虑下列性质:}
\[
  (+_{\!f})\qquad
  \text{对每个 \(s\in S\), \(T_s\) 可由
  \(G_s\) 的闭子概形表出.}
\]
\emph{设 \(u,w:X\to G\) 和 \(v:Y\to G\) 为
\(S\)-概形的态射. 则:}
\begin{enumerate}
\item[\textup{(i)}]
  \(\underline{\operatorname{Transp}}(u,w)\) 和
  \(\underline{\operatorname{Centr}}(u)
    =\underline{\operatorname{Transp}}(u,u)\)
  满足条件 \((+_{\!f})\).
\item[\textup{(ii)}]
  若对每个 \(s\in S\), \(v_s\) 为闭
  浸入, 则 \(\underline{\operatorname{Transp}}_G(u(X),v(Y))\) 和
  \(\underline{\operatorname{Norm}}_G v(Y)\) 满足条件
  \((+_{\!f})\).
\item[\textup{\SGAInlineRefTarget{prop:sga3-VIB-6.2:item:iii}(iii)}]
  若对每个 \(s\in S\),
  \(u_s\) 和 \(v_s\) 为闭浸入, 则
  \(\underline{\operatorname{Transp}}^{\mathrm{str}}_G(X,Y)\)
  满足条件 \((+_{\!f})\).
\end{enumerate}

这由注~\hyperref[rem:sga3-VIB-6.1.1]{6.1.1}
和下面的推论~\hyperref[cor:sga3-VIB-6.2.5]{6.2.5}
得到.%
\footnote{\label{note:sga3-VIB-52}编者注: 原文
引用了第 VIII 篇 \S6 的结果. 为便于读者查阅,
我们在下面的 \SGARefLink{def:sga3-VIB-6.2.1}{6.2.1}--\SGARefLink{cor:sga3-VIB-6.2.5}{6.2.5}
中重述了这些结果 (除 VIII, \SGARefLink{sga3:viii:proposition:6:3}{6.3} 和 \SGARefLink{sga3:viii:remark:6:8}{6.8} 以外).
此外, A.\ Grothendieck 在 VIII.6 开头的一则注中已经
提出这一建议: ``本节独立于可对角化群理论;
它自然应当置于
VI\(_B\) 中.''}

\medskip
\phantomsection\label{def:sga3-VIB-6.2.1}%
\noindent\textbf{定义 6.2.1.---}
设 \(f:X\to S\) 为概形态射. 若能找到 \(S\) 的一个
仿射开覆盖 \(S_i\), 对每个 \(i\) 找到一个在 \(S_i\) 上
忠实平坦的仿射 \(S_i\)-概形 \(S'_i\), 并找到
\[
  X'_i=X\times_S S'_i
\]
的一个由仿射开子集 \(X'_{ij}\) 组成的覆盖 \((X'_{ij})_j\),
使得对每个 \((i,j)\), \(X'_{ij}\) 的环是 \(S'_i\) 的环上的投射模,
则称 \(f\) \emph{本质自由 (essentially free)}, 或称 \(X\)
\emph{在 \(S\) 上本质自由}.%
\footnote{\label{note:sga3-VIB-53}编者注: 一方面,
我们按 VIII, \SGARefLink{sga3:viii:remark:6:8}{注 6.8} 中的说明,
将 ``自由'' 替换为 ``投射''. 另一方面, 本质自由
\(S\)-概形这一概念与 M.\ Raynaud 和 L.\ Gruson
引入并研究的平坦且纯的 \(S\)-概形这一几何概念
密切相关 ([RG71]); 见下面增补的
\hyperref[par:sga3-VIB-6.9]{6.9}.}

\medskip
\phantomsection\label{prop:sga3-VIB-6.2.2}%
\noindent\textbf{命题 6.2.2.---}
\begin{enumerate}
\item[\textup{a)}]
  \emph{若 \(X\) 在 \(S\) 上本质自由, 则它在
  \(S\) 上平坦; 当 \(S\) 为 Artin 概形时, 反之也成立.}
\item[\textup{b)}]
  \emph{若 \(S\) 为一个域的谱, 则每个 \(S\)-概形
  都在 \(S\) 上本质自由.}
\item[\textup{c)}]
  \emph{若 \(X\) 在 \(S\) 上本质自由, 则对每个
  \(S'\to S\), \(X'=X\times_S S'\) 在 \(S'\) 上本质自由.
  若 \(S'\to S\) 忠实平坦且拟紧, 则反之也成立.}
\end{enumerate}

证明是直接的; 对于 \textup{a)} 的逆命题,
使用 Artin 局部环上的平坦模为自由模这一事实.%
\footnote{\label{note:sga3-VIB-54}编者注: 事实上, 设
\((A,\mathfrak m)\) 为 Artin 局部环, \(k\) 为其剩余
域, \(M\) 为任意 \(A\)-模, 并设
\((x_i)_{i\in I}\) 为 \(M\) 中的一族元素, 其像构成
\(M/\mathfrak mM\) 在 \(k\) 上的一组基. 设 \(F\) 为以 \((e_i)_{i\in I}\)
为基的自由 \(A\)-模, 并设
\(\varphi:F\to M\) 为由
\(\varphi(e_i)=x_i\) 定义的 \(A\)-态射. 则 \(Q=\operatorname{Coker}\varphi\)
满足 \(Q=\mathfrak mQ\), 从而由于 \(\mathfrak m\) 幂零,
\(Q=0\). 若再假设 \(M\) 在 \(A\) 上平坦,
则 \(K=\operatorname{Ker}\varphi\) 满足
\(K\otimes_Ak=0\), 即 \(K=\mathfrak mK\), 故 \(K=0\).}

引入定义~\hyperref[def:sga3-VIB-6.2.1]{6.2.1}
的理由在于下述定理.

% 翻译统计: 源行数 111; 输出行数 110 (含本统计行); 行内/行间数学片段 77; 翻译占位 0. 底本疑点另报主线程.
